\section{Ajuste de hiperparámetros y reproducción de experimentos}

\subsection{Combinaciones comunes de hiperparámetros}

En los experimentos de CS224R, Chelsea y su equipo obtuvieron un desempeño estable con las siguientes combinaciones:

\begin{itemize}
    \item \textbf{PPO}: clip $\epsilon=0.1$, tamaño de lote=64, learning rate=2e-4, value coeff=0.5, entropy coeff=0.01.
    \item \textbf{SAC}: Q learning rate=3e-4, policy learning rate=3e-4, replay buffer size=10e5, tamaño de lote=256, target smoothing coeff=0.005.
    \item \textbf{Replay buffer}: prioritized sampling weight $\alpha=0.6$, importance sampling exponent $\beta$ crece linealmente de 0.4 a 1.0.
\end{itemize}

\subsection{Puntos clave de la reproducción}

Reproducir experimentos de forma adversarial requiere fijar la semilla aleatoria y utilizar un deterministic rollout. Se recomienda registrar en cada run

\begin{itemize}
    \item seed, env name, max episode length
    \item optimizer state (Adam momenta) y learning rate schedule
    \item buffer size, prioritization hyperparameters
\end{itemize}

\begin{warningbox}{Sensibilidad a la semilla}
Para prevenir fallas esporádicas, es necesario incluir en el logging el return, la length y la std de las advantages de cada rollout. Si la return variance supera lo esperado, se debe detener de inmediato el entrenamiento e investigar si el replay buffer está saturado de high-loss transitions.
\end{warningbox}

\subsection{Resumen de la sección}

La combinación de hiperparámetros y el registro de logs son fundamentales para la reproducción de experimentos. Fijar el seed, registrar el optimizer state, monitorear la variance y establecer un umbral de interrupción automática puede reducir de manera significativa la probabilidad de long-tail failure.
